\clearpage\thispagestyle{empty}
\pdfbookmark[0]{Quellen zur Besprechung von Cagnolis Sternverzeichnis}{cagnoli1808-sources}
\begin{center}\textbf{Quellen zur Besprechung von Cagnolis Sternverzeichnis}\end{center}
\noindent Diese KI-Übertragung ins heutige Deutsch folgt der Besprechung von Gauß auf S.~496--497 der \href{https://gdz.sub.uni-goettingen.de/id/PPN236007467}{\textit{Werke}, Bd.~VI} (Göttingen, 1874). Sie erschien ohne Namensnennung unter der Überschrift \emph{Modena} in den \textit{Göttingischen gelehrten Anzeigen}, Stück~77, 14. Mai 1808, S.~\href{https://api.digitale-sammlungen.de/iiif/image/v2/bsb10538673_00767/full/1800,/0/default.jpg}{761}--\href{https://api.digitale-sammlungen.de/iiif/image/v2/bsb10538673_00769/full/1800,/0/default.jpg}{763}. Das verlinkte Exemplar stammt aus der Bayerischen Staatsbibliothek.
\par\smallskip\noindent Der französische Titel bedeutet \emph{Verzeichnis von 501 Sternen mit nachfolgenden Tafeln zur Aberration und Nutation}, von Antoine Cagnoli. Gauß beschreibt das Verzeichnis und die Tafeln. KI-Übertragung und KI-Satz.
\clearpage
